\documentclass[12pt,a4paper]{article}
\usepackage[spanish]{babel} % Corta palabras en espa~nol
\usepackage[utf8]{inputenc} % Escribir con acentos, ~n...
\usepackage{eurosym} % smbolo del euro
\begin{document}
%\secttion{Introducción}
\leftline{Introducción}
\medskip
Este primer ejemplo trata de demostrar la facilidad de
\LaTeX{}. Por ejemplo varios espacios en blanco
se tratan como uno.
Para empezar un nuevo párrafo basta dejar una llínea en blanco.  Expresiones matemáticas son sencillas de
escribir \footnote{nota al pie}:
$a=\sum_{i=1}^{i=\infty} x_i^{n+1}$ y deben ser escritas
entre dólares. Los superíndices se obtienen con \^{ },
$x^3 y^{\alpha + \beta}$, mientras que los suníndices son con \_.  Además se puede escribir la
fórmula centrada
\[ z^{2+nalpha}_{n+k}. \]
\medskip
El símbolo del euro \euro{} existe.
\end{document}